
\documentclass[12pt,a4paper]{article}

% ------------------------------------------------
% Paquetes
% ------------------------------------------------
\usepackage[utf8]{inputenc}
\usepackage[T1]{fontenc}
\usepackage[spanish, es-nodecimaldot]{babel}
\usepackage{geometry}
\usepackage{graphicx}
\usepackage{float}
\usepackage{xcolor}
\usepackage{listings}
\usepackage{hyperref}
\usepackage{fancyhdr}
\usepackage{titlesec}

% ------------------------------------------------
% Configuración de página
% ------------------------------------------------
\geometry{
    top=2.5cm,
    bottom=2.5cm,
    left=2.5cm,
    right=2.5cm
}

% ------------------------------------------------
% Configuración de listings
% ------------------------------------------------
\definecolor{codegray}{gray}{0.95}

\lstset{
    backgroundcolor=\color{codegray},
    basicstyle=\ttfamily\small,
    breaklines=true,
    frame=single,
    columns=fullflexible,
    keepspaces=true,
    showstringspaces=false,
    tabsize=2
}

% ------------------------------------------------
% Encabezado y pie de página
% ------------------------------------------------
\pagestyle{fancy}
\fancyhf{}
\lhead{Interacción de redes}
\rhead{Práctica 7}
\cfoot{\thepage}
\setlength{\headheight}{15pt}

% ------------------------------------------------
% Configuración de títulos
% ------------------------------------------------
\titleformat{\section}
{\large\bfseries}
{\thesection.}
{0.5em}
{}

% ------------------------------------------------
% Hipervínculos
% ------------------------------------------------
\hypersetup{
    colorlinks=true,
    linkcolor=black,
    urlcolor=blue,
    citecolor=black
}

% ------------------------------------------------
% Comando para insertar capturas opcionales
% ------------------------------------------------
\newcommand{\captura}[2]{%
    \begin{figure}[H]
        \centering
        \IfFileExists{#1}{%
            \includegraphics[
                width=0.92\textwidth,
                keepaspectratio
            ]{#1}
        }{%
            \fbox{%
                \parbox[c][4cm][c]{0.82\textwidth}{%
                    \centering
                    \textbf{Captura pendiente}\\[0.5em]
                    Archivo esperado: \texttt{#1}
                }%
            }%
        }
        \caption{#2}
    \end{figure}
}

% ------------------------------------------------
% Datos del documento
% ------------------------------------------------
\title{
    Práctica 7:
    \textbf{Configuración de ACL Extendidas
    en Equipo Físico}\\[0.5em]
    \large Interacción de redes
}

\author{
    Jorge Alberto Suárez Saldaña\\
    \small Matrícula: 355992\\
    \small Profesor:
    Raymundo Antonio González Grimaldo
}

\date{\today}

% ------------------------------------------------
% Documento
% ------------------------------------------------
\begin{document}

\maketitle

\section{Introducción}

Las listas de control de acceso, conocidas como ACL
(\textit{Access Control Lists}), son mecanismos utilizados
en dispositivos de red para permitir o denegar el
tráfico de acuerdo con criterios previamente establecidos.

En los routers Cisco, las ACL extendidas permiten
evaluar características como las direcciones IPv4
de origen y destino, el protocolo de transporte,
los puertos utilizados y determinados tipos
de mensajes ICMP.

A diferencia de las ACL estándar, que filtran
principalmente según la dirección IPv4 de origen,
las ACL extendidas permiten establecer políticas
de seguridad más específicas.

En esta práctica se implementó una topología de
red física compuesta por dos routers Cisco,
dos switches, dos equipos cliente y dos servidores
virtualizados con Ubuntu Server.

Se configuró el protocolo OSPF para permitir
la comunicación entre las diferentes subredes.
Posteriormente, se implementaron ACL extendidas
para restringir el acceso al servidor de bases
de datos y controlar las solicitudes ICMP entre
las redes interna y externa.

Adicionalmente, se desarrolló una aplicación
web utilizando Apache y PHP, capaz de consultar
información almacenada en un servidor remoto
MySQL/MariaDB.

Esta aplicación permitió comprobar que el
filtrado de tráfico no interfiriera con
las comunicaciones autorizadas.

\section{Topología de red}

La topología implementada comprende dos
redes locales comunicadas mediante un enlace serial.

El router RI proporciona conectividad a
la red interna, mientras que el router RE
proporciona conectividad a la red externa.

Las redes utilizadas fueron:

\begin{itemize}
    \item Red interna:
    \texttt{52.20.30.0/24}.

    \item Red externa:
    \texttt{140.10.20.0/24}.

    \item Enlace serial:
    \texttt{200.0.0.0/30}.
\end{itemize}

El switch SI conecta al router RI
con PCI y Web Server.

El switch SE conecta al router RE
con PCE y DB Server.

Los servidores se ejecutaron en máquinas
virtuales, conectadas mediante adaptadores
de red en modo puente
(\textit{Bridged Networking}).

Este modo permitió que las máquinas virtuales
participaran directamente en las redes
del laboratorio.

\captura{topologia.png}{
    Topología física utilizada en la práctica.
}

\subsection{Tabla de direccionamiento}

\begin{table}[H]
\centering
\small
\begin{tabular}{|l|l|l|l|}
\hline
\textbf{Dispositivo} &
\textbf{Interfaz} &
\textbf{IPv4} &
\textbf{Máscara} \\
\hline
RI & G0/0 & 52.20.30.1 & /24 \\
\hline
RI & Serial & 200.0.0.1 & /30 \\
\hline
RE & Serial & 200.0.0.2 & /30 \\
\hline
RE & G0/0 & 140.10.20.1 & /24 \\
\hline
Web Server & Ethernet & 52.20.30.10 & /24 \\
\hline
PCI & Ethernet & 52.20.30.20 & /24 \\
\hline
DB Server & Ethernet & 140.10.20.10 & /24 \\
\hline
PCE & Ethernet & 140.10.20.20 & /24 \\
\hline
\end{tabular}
\caption{Direccionamiento IPv4 de la topología.}
\end{table}

Las interfaces Ethernet de los routers
utilizan una máscara \texttt{/24},
equivalente a \texttt{255.255.255.0}.

El enlace serial utiliza una máscara
\texttt{/30}, equivalente a
\texttt{255.255.255.252}.

Las puertas de enlace de los equipos
finales se configuraron de la
siguiente manera:

\begin{itemize}
    \item Web Server y PCI:
    \texttt{52.20.30.1}.

    \item DB Server y PCE:
    \texttt{140.10.20.1}.
\end{itemize}

\section{Preparación del servidor web}

\subsection{Instalación del sistema operativo}

Para Web Server se utilizó una
máquina virtual con Ubuntu Server
26.04.1 LTS, denominada
\texttt{webserverpract}.

La versión del sistema operativo
se verificó utilizando:

\begin{lstlisting}
cat /etc/os-release
\end{lstlisting}

El sistema reportó:

\begin{lstlisting}
PRETTY_NAME="Ubuntu 26.04.1 LTS"
VERSION_ID="26.04"
VERSION_CODENAME=resolute
\end{lstlisting}

Inicialmente se intentó instalar PHP 8.3,
pero los repositorios configurados
no proporcionaban dichos paquetes.

Por esta razón se utilizó PHP 8.5,
disponible para la versión instalada
de Ubuntu.

Este cambio no modificó los objetivos
de comunicación y filtrado de tráfico
de la práctica.

\subsection{Instalación de Apache y PHP}

Primero se actualizaron los índices
de paquetes:

\begin{lstlisting}
sudo apt update
\end{lstlisting}

Posteriormente se instalaron los servicios
necesarios:

\begin{lstlisting}
sudo apt install -y \
    apache2 \
    php8.5 \
    php8.5-cli \
    libapache2-mod-php8.5 \
    php8.5-mysql \
    mariadb-client \
    curl \
    net-tools
\end{lstlisting}

El paquete \texttt{php8.5-mysql}
proporciona los controladores necesarios
para establecer conexiones entre PHP
y MySQL/MariaDB.

\subsection{Activación y verificación de Apache}

Se habilitó Apache para iniciar
automáticamente con el sistema:

\begin{lstlisting}
sudo systemctl enable --now apache2
\end{lstlisting}

Se comprobó su estado mediante:

\begin{lstlisting}
sudo systemctl status apache2
\end{lstlisting}

El servicio reportó:

\begin{lstlisting}
Active: active (running)
\end{lstlisting}

Se verificó también la versión de PHP:

\begin{lstlisting}
php -v
\end{lstlisting}

La versión instalada fue PHP 8.5.4.

\subsection{Prueba de ejecución de PHP}

Para comprobar la integración entre
Apache y PHP se creó un archivo temporal:

\begin{lstlisting}
echo '<?php echo "Apache y PHP funcionando. Version: " . PHP_VERSION; ?>' \
    | sudo tee /var/www/html/check.php
\end{lstlisting}

La prueba se realizó mediante:

\begin{lstlisting}
curl http://localhost/check.php
\end{lstlisting}

La respuesta obtenida fue:

\begin{lstlisting}
Apache y PHP funcionando. Version: 8.5.4
\end{lstlisting}

Este resultado confirmó que Apache
interpretaba correctamente el código PHP.

Finalmente se eliminó el archivo temporal:

\begin{lstlisting}
sudo rm /var/www/html/check.php
\end{lstlisting}

\section{Configuración del servidor de bases de datos}

El segundo servidor se destinó al
almacenamiento y administración
de una base de datos relacional.

El gestor utilizado fue compatible
con el protocolo MySQL, permitiendo
la conexión desde PHP mediante PDO.

\subsection{Base de datos y estructura}

La base de datos utilizada recibió
el nombre:

\begin{lstlisting}
laboratorio_acl
\end{lstlisting}

Dentro de ella se creó una tabla
denominada \texttt{productos}.

Su estructura fue la siguiente:

\begin{table}[H]
\centering
\begin{tabular}{|l|l|}
\hline
\textbf{Columna} & \textbf{Tipo} \\
\hline
id & INT \\
\hline
nombre & VARCHAR(100) \\
\hline
precio & DECIMAL(10,2) \\
\hline
cantidad & INT \\
\hline
fecha\_registro & DATE \\
\hline
disponible & TINYINT(1) \\
\hline
\end{tabular}
\caption{Estructura de la tabla productos.}
\end{table}

La tabla contiene seis columnas
con diferentes tipos de datos,
cumpliendo el requisito establecido
en la práctica.

Los campos representan:

\begin{itemize}
    \item \texttt{id}: identificador
    del producto.

    \item \texttt{nombre}: nombre
    del artículo.

    \item \texttt{precio}: precio
    del producto.

    \item \texttt{cantidad}: unidades
    registradas.

    \item \texttt{fecha\_registro}:
    fecha de incorporación.

    \item \texttt{disponible}:
    indicador booleano de disponibilidad.
\end{itemize}

\subsection{Consultas de verificación}

Para inspeccionar la base de datos
se utilizaron las instrucciones:

\begin{lstlisting}[language=SQL]
SHOW DATABASES;

USE laboratorio_acl;

SHOW TABLES;

DESCRIBE productos;

SELECT * FROM productos;
\end{lstlisting}


\subsection{Conexiones remotas}

Para permitir la comunicación
entre Web Server y DB Server,
el gestor debía aceptar conexiones
TCP en el puerto 3306.

En la topología definitiva,
el servidor de bases de datos
utilizó la dirección:

\begin{lstlisting}
140.10.20.10
\end{lstlisting}

La configuración del servicio
debía permitir conexiones sobre
esa interfaz, mediante un parámetro
equivalente a:

\begin{lstlisting}
bind-address = 140.10.20.10
\end{lstlisting}

El servicio se verificó mediante:

\begin{lstlisting}
sudo ss -lntp | grep 3306
\end{lstlisting}

Se utilizó el usuario
\texttt{webuser} para las consultas
desde la aplicación PHP.

Como medida de seguridad,
se recomienda otorgar exclusivamente
el permiso \texttt{SELECT} a la
cuenta utilizada por la aplicación.

Una configuración de permisos
apropiada es:

\begin{lstlisting}[language=SQL]
CREATE USER IF NOT EXISTS
'webuser'@'52.20.30.10'
IDENTIFIED BY 'CLAVE_PRIVADA';

GRANT SELECT
ON laboratorio_acl.productos
TO 'webuser'@'52.20.30.10';
\end{lstlisting}

La contraseña real no se incluye
en este reporte por motivos
de seguridad.

\section{Desarrollo de la aplicación web}

\subsection{Arquitectura de comunicación}

La aplicación se diseñó utilizando
una arquitectura cliente-servidor.

El flujo de información fue:

\begin{enumerate}
    \item Un cliente realiza una
    solicitud HTTP a Web Server.

    \item Apache recibe la solicitud
    y ejecuta el archivo PHP.

    \item PHP establece una conexión
    TCP con DB Server.

    \item El gestor ejecuta una consulta SQL.

    \item PHP recibe los registros
    y genera una tabla HTML.

    \item Apache devuelve la página
    al navegador del cliente.
\end{enumerate}

Los servicios utilizan los siguientes
puertos:

\begin{itemize}
    \item HTTP: TCP 80.
    \item MySQL/MariaDB: TCP 3306.
\end{itemize}

\subsection{Archivo de conexión}

Se creó un archivo de configuración
fuera del directorio público de Apache:

\begin{lstlisting}
sudo install -d -m 750 \
    -o root -g www-data \
    /etc/practica-web

sudo nano /etc/practica-web/db.php
\end{lstlisting}

El archivo contiene los parámetros
necesarios para la conexión:

\begin{lstlisting}
<?php

return [
    'host' => '140.10.20.10',
    'database' => 'laboratorio_acl',
    'username' => 'webuser',
    'password' => 'CLAVE_PRIVADA'
];
\end{lstlisting}

Se restringieron sus permisos:

\begin{lstlisting}
sudo chown root:www-data \
    /etc/practica-web/db.php

sudo chmod 640 \
    /etc/practica-web/db.php
\end{lstlisting}

La separación de este archivo
permite modificar los parámetros
de conexión sin alterar
la lógica de la aplicación.

\subsection{Código PHP implementado}

La aplicación principal se almacenó en:

\begin{lstlisting}
/var/www/html/index.php
\end{lstlisting}

La lógica PHP utilizada para
consultar la base de datos fue:

\begin{lstlisting}
<?php

$config = require '/etc/practica-web/db.php';

$registros = [];
$error = null;

function escapar($valor): string {
    return htmlspecialchars(
        (string) $valor,
        ENT_QUOTES | ENT_SUBSTITUTE,
        'UTF-8'
    );
}

try {
    $dsn = "mysql:host={$config['host']};" .
           "port=3306;" .
           "dbname={$config['database']};" .
           "charset=utf8mb4";

    $pdo = new PDO(
        $dsn,
        $config['username'],
        $config['password'],
        [
            PDO::ATTR_ERRMODE =>
                PDO::ERRMODE_EXCEPTION,
            PDO::ATTR_DEFAULT_FETCH_MODE =>
                PDO::FETCH_ASSOC,
            PDO::ATTR_TIMEOUT => 5
        ]
    );

    $sql = "SELECT id, nombre, precio,
                   cantidad, fecha_registro,
                   disponible
            FROM productos
            ORDER BY id";

    $registros =
        $pdo->query($sql)->fetchAll();

} catch (Throwable $e) {
    error_log(
        "Error de base de datos: " .
        $e->getMessage()
    );

    $error =
        "No fue posible consultar la base de datos.";

    http_response_code(503);
}

?>
\end{lstlisting}

En el documento HTML se utilizó
un ciclo \texttt{foreach} para
mostrar los registros obtenidos:

\begin{lstlisting}
<table>
<thead>
<tr>
    <th>ID</th>
    <th>Nombre</th>
    <th>Precio</th>
    <th>Cantidad</th>
    <th>Fecha de registro</th>
    <th>Disponible</th>
</tr>
</thead>

<tbody>
<?php foreach ($registros as $fila): ?>
<tr>
    <td><?= escapar($fila['id']) ?></td>
    <td><?= escapar($fila['nombre']) ?></td>
    <td><?= escapar($fila['precio']) ?></td>
    <td><?= escapar($fila['cantidad']) ?></td>
    <td><?= escapar($fila['fecha_registro']) ?></td>
    <td>
        <?= $fila['disponible'] ? 'Si' : 'No' ?>
    </td>
</tr>
<?php endforeach; ?>
</tbody>
</table>
\end{lstlisting}

El código final también incorporó
estilos CSS para presentar los
registros en una tabla legible.

La función \texttt{htmlspecialchars}
se utilizó para evitar que los
datos recibidos de la base
se interpretaran como código HTML.

Asimismo, se implementó un mecanismo
de manejo de excepciones que devuelve
el estado HTTP 503 si la aplicación
no consigue consultar la base de datos.

\subsection{Verificación del código}

La sintaxis de la aplicación se
comprobó con:

\begin{lstlisting}
php -l /var/www/html/index.php
\end{lstlisting}

Se obtuvo:

\begin{lstlisting}
No syntax errors detected
in /var/www/html/index.php
\end{lstlisting}

Después se reinició Apache:

\begin{lstlisting}
sudo systemctl restart apache2
\end{lstlisting}

Se eliminó la página predeterminada:

\begin{lstlisting}
sudo rm -f /var/www/html/index.html
\end{lstlisting}

Finalmente, se utilizó:

\begin{lstlisting}
curl -i http://localhost/
\end{lstlisting}

Antes de establecer la conexión
con el servidor de bases de datos,
la aplicación devolvió
\texttt{503 Service Unavailable}.

Una vez configurada la conexión
remota, la aplicación mostró
correctamente los registros.

\captura{paginaWeb.png}{
    Aplicación PHP mostrando los productos
    consultados desde el servidor remoto.
}

\section{Pruebas iniciales en red provisional}

Antes de disponer de los routers
físicos, se realizaron pruebas
entre ambas máquinas virtuales
en una red local alternativa.

Las direcciones utilizadas fueron:

\begin{table}[H]
\centering
\begin{tabular}{|l|l|}
\hline
\textbf{Dispositivo} & \textbf{IPv4} \\
\hline
Web Server & 10.103.150.187 \\
\hline
DB Server & 10.103.150.181 \\
\hline
\end{tabular}
\caption{Direccionamiento provisional
para las pruebas de aplicación.}
\end{table}

El archivo de configuración PHP
se ajustó temporalmente para
conectarse a \texttt{10.103.150.181}.

Se efectuaron pruebas de conectividad
y del puerto MySQL:

\begin{lstlisting}
ping -c 4 10.103.150.181

nc -vz -w 5 10.103.150.181 3306
\end{lstlisting}

La conexión SQL se comprobó mediante:

\begin{lstlisting}
mariadb \
    -h 10.103.150.181 \
    -u webuser \
    -p \
    laboratorio_acl
\end{lstlisting}

Posteriormente se consultó:

\begin{lstlisting}[language=SQL]
SELECT * FROM productos;
\end{lstlisting}

Las pruebas confirmaron que
la aplicación podía comunicarse
con el servidor de bases
de datos a través de la red.

Se comprobó además que la
página web mostraba los registros
correctamente.

Esta etapa permitió separar los
problemas de instalación y
configuración de servicios
de los problemas relacionados
con el enrutamiento físico.

\section{Configuración de OSPF}

Una vez implementada la topología
en el laboratorio físico, se utilizó
OSPF para comunicar las redes
interna y externa.

Ambos routers participaron en
el área 0 del proceso OSPF 10.

\subsection{Configuración del router RI}

La configuración de referencia
utilizada para las redes fue:

\begin{lstlisting}
enable
configure terminal

hostname RI

interface GigabitEthernet0/0
 ip address 52.20.30.1 255.255.255.0
 no shutdown
 exit

interface Serial0/0/0
 ip address 200.0.0.1 255.255.255.252
 no shutdown
 exit

router ospf 10
 router-id 1.1.1.1
 passive-interface GigabitEthernet0/0
 network 52.20.30.0 0.0.0.255 area 0
 network 200.0.0.0 0.0.0.3 area 0
 exit

end
\end{lstlisting}

\subsection{Configuración del router RE}

\begin{lstlisting}
enable
configure terminal

hostname RE

interface GigabitEthernet0/0
 ip address 140.10.20.1 255.255.255.0
 no shutdown
 exit

interface Serial0/0/0
 ip address 200.0.0.2 255.255.255.252
 no shutdown
 exit

router ospf 10
 router-id 2.2.2.2
 passive-interface GigabitEthernet0/0
 network 140.10.20.0 0.0.0.255 area 0
 network 200.0.0.0 0.0.0.3 area 0
 exit

end
\end{lstlisting}

Los nombres de interfaz
corresponden al esquema de
configuración propuesto.
En equipo físico deben coincidir
con las interfaces instaladas.

\subsection{Verificación del enrutamiento}

Para comprobar OSPF se utilizaron
los siguientes comandos:

\begin{lstlisting}
show ip interface brief
show ip ospf neighbor
show ip route
show ip route ospf
\end{lstlisting}

El resultado esperado consistía
en establecer una vecindad OSPF
en estado \texttt{FULL} y
aprender las redes remotas.

RI debe aprender:

\begin{lstlisting}
O 140.10.20.0/24
\end{lstlisting}

RE debe aprender:

\begin{lstlisting}
O 52.20.30.0/24
\end{lstlisting}

La comunicación posterior
entre los servidores ubicados
en diferentes redes confirmó
que existía un mecanismo de
enrutamiento funcional.

\section{Configuración definitiva de los servidores}

Una vez disponibles las interfaces
y los switches del laboratorio,
se modificó el direccionamiento
de los servidores.

Web Server utilizó:

\begin{lstlisting}
IP:       52.20.30.10
Mascara:  255.255.255.0
Gateway:  52.20.30.1
\end{lstlisting}

DB Server utilizó:

\begin{lstlisting}
IP:       140.10.20.10
Mascara:  255.255.255.0
Gateway:  140.10.20.1
\end{lstlisting}

La configuración de Ubuntu
se realizó mediante Netplan.

Para Web Server se utilizó
una configuración equivalente a:

\begin{lstlisting}
network:
  version: 2
  ethernets:
    enp0s1:
      dhcp4: false
      addresses:
        - 52.20.30.10/24
      routes:
        - to: default
          via: 52.20.30.1
\end{lstlisting}

El identificador \texttt{enp0s1}
es ilustrativo y depende de
la interfaz de cada máquina virtual.

Los cambios se validaron mediante:

\begin{lstlisting}
sudo netplan generate
sudo netplan try
\end{lstlisting}

Posteriormente se comprobó el
direccionamiento y la tabla
de rutas con:

\begin{lstlisting}
ip -br addr
ip route
\end{lstlisting}

También se actualizaron
los parámetros de conexión
de PHP, así como la dirección
de escucha y los permisos del
servidor MySQL/MariaDB.


\section{Pruebas de conectividad previas a las ACL}

Antes de aplicar restricciones
se comprobó la comunicación
entre las redes.

Desde Web Server se utilizaron
comandos como:

\begin{lstlisting}
ping -c 4 52.20.30.1
ping -c 4 140.10.20.1
ping -c 4 140.10.20.10
\end{lstlisting}

También se verificó la conexión
al servidor de bases de datos:

\begin{lstlisting}
mariadb \
    -h 140.10.20.10 \
    -u webuser \
    -p \
    laboratorio_acl
\end{lstlisting}

Dentro del gestor se ejecutó:

\begin{lstlisting}[language=SQL]
SELECT * FROM productos;
\end{lstlisting}

Finalmente, se comprobó la
aplicación mediante:

\begin{lstlisting}
curl -i http://localhost/
\end{lstlisting}

La página permitió visualizar
los productos almacenados en
DB Server.

El funcionamiento de esta
aplicación sobre la topología
física confirmó la comunicación
entre los servidores.

Adicionalmente, se contemplaron
pruebas ICMP entre PCI y PCE,
así como entre las estaciones
y los servidores.

Estas pruebas son necesarias
para distinguir los bloqueos
intencionales de las ACL de
posibles fallos de conectividad.

\captura{pings.png}{
    Pruebas de conectividad
    realizadas antes de aplicar ACL.
}

\section{Políticas de seguridad}

Las políticas establecidas
para esta práctica fueron:

\begin{enumerate}
    \item Proteger el acceso a DB Server:
    únicamente Web Server debe poder
    establecer comunicación con él.

    \item Permitir las solicitudes
    de ping desde
    \texttt{52.20.30.0/24}
    hacia
    \texttt{140.10.20.0/24}.

    \item Impedir las solicitudes
    de ping desde
    \texttt{140.10.20.0/24}
    hacia
    \texttt{52.20.30.0/24}.
\end{enumerate}

Para implementar estas políticas
se utilizaron dos ACL extendidas
sobre el router RE.

Las ACL recibieron los nombres:

\begin{itemize}
    \item \texttt{BLOQUEAR\_PING\_EXTERNO}.
    \item \texttt{PROTEGER\_DB}.
\end{itemize}

Ambas se aplicaron sobre
la interfaz que conecta
RE con el switch SE,
pero en direcciones distintas.

\section{Configuración de ACL extendidas}

\subsection{ACL para bloquear ping desde la red externa}

La primera ACL impide que
los equipos de la red externa
inicien solicitudes ICMP
hacia la red interna.

En RE se configuró:

\begin{lstlisting}
enable
configure terminal

ip access-list extended BLOQUEAR_PING_EXTERNO
 remark Bloquear ping de red externa a interna
 deny icmp 140.10.20.0 0.0.0.255 52.20.30.0 0.0.0.255 echo
 remark Permitir respuestas y otros protocolos
 permit ip any any
 exit
\end{lstlisting}

La instrucción \texttt{deny icmp}
filtra únicamente los mensajes
de tipo \texttt{echo} que
coincidan con las direcciones
de origen y destino indicadas.

Los mensajes \texttt{echo-reply}
no se bloquean mediante esa regla.

Esto permite que un equipo
de la red interna inicie un ping
hacia la red externa y reciba
su respuesta.

\subsection{ACL para proteger DB Server}

La segunda ACL permite que
Web Server se comunique con
DB Server y deniega las
comunicaciones de otros
orígenes que atraviesen RE.

\begin{lstlisting}
ip access-list extended PROTEGER_DB
 remark Permitir Web Server hacia DB Server
 permit ip host 52.20.30.10 host 140.10.20.10

 remark Denegar otros origenes hacia DB
 deny ip any host 140.10.20.10

 remark Permitir el resto del trafico
 permit ip any any
 exit
\end{lstlisting}

La primera regla:

\begin{lstlisting}
permit ip host 52.20.30.10 host 140.10.20.10
\end{lstlisting}

permite el tráfico IPv4
originado específicamente
por Web Server y dirigido
a DB Server.

La segunda regla:

\begin{lstlisting}
deny ip any host 140.10.20.10
\end{lstlisting}

bloquea el resto del tráfico
con destino a DB Server
cuando dicho tráfico
atraviesa la interfaz
controlada por la ACL.

La instrucción final:

\begin{lstlisting}
permit ip any any
\end{lstlisting}

permite el tráfico que no
coincide con las reglas anteriores.

Sin esta instrucción,
el \textit{deny} implícito
de la ACL podría bloquear
otras comunicaciones legítimas.

\subsection{Aplicación de las ACL}

Las dos ACL se asociaron
a la interfaz LAN de RE.

\begin{lstlisting}
interface GigabitEthernet0/0
 ip access-group BLOQUEAR_PING_EXTERNO in
 ip access-group PROTEGER_DB out
 exit

end
\end{lstlisting}

La primera se aplica
en dirección de entrada,
mientras que la segunda
se aplica en dirección
de salida.

La ACL de entrada analiza
el tráfico recibido desde SE.

La ACL de salida analiza
el tráfico que RE entrega
hacia la red externa.

Finalmente, se guardó
la configuración:

\begin{lstlisting}
copy running-config startup-config
\end{lstlisting}

\subsection{Verificación de las ACL}

Para revisar las reglas
se utilizaron los comandos:

\begin{lstlisting}
show ip access-lists
show ip interface GigabitEthernet0/0
\end{lstlisting}

El resultado esperado incluye:

\begin{lstlisting}
Inbound access list is BLOQUEAR_PING_EXTERNO
Outgoing access list is PROTEGER_DB
\end{lstlisting}

El comando \texttt{show ip access-lists}
permite observar las reglas
y sus contadores de coincidencias.

Estos contadores ayudan
a identificar qué entradas
están permitiendo o denegando
el tráfico.

\section{Consideración sobre el tráfico local}

Una observación importante
de esta topología es que
PCE y DB Server se encuentran
en la misma subred:

\begin{lstlisting}
140.10.20.0/24
\end{lstlisting}

Ambos dispositivos están
conectados al switch SE.

Por lo tanto, las comunicaciones
directas entre estos equipos
no requieren atravesar
el router RE.

Esto significa que
las ACL configuradas
en RE no pueden filtrar
por sí solas el tráfico
local entre PCE y DB Server.

Para implementar completamente
la política de protección
del servidor se requiere
un mecanismo adicional.

Una alternativa consiste
en configurar reglas de firewall
en Ubuntu Server para permitir
únicamente conexiones autorizadas
desde Web Server.

Otra alternativa consiste
en modificar la topología
y separar PCE y DB Server
en diferentes VLAN o subredes,
obligando a enrutar su tráfico.

Esta limitación debe
considerarse al evaluar
el cumplimiento de la
primera política de seguridad.

La configuración de las ACL
en RE protege las comunicaciones
enrutadas hacia DB Server,
pero no sustituye el filtrado
local en la misma VLAN.

\section{Pruebas posteriores a las ACL}

Después de aplicar las ACL,
las pruebas deben comprobar
que los flujos autorizados
continúan funcionando y
que los flujos restringidos
son bloqueados.

Los resultados esperados
según las políticas son:

\begin{table}[H]
\centering
\small
\begin{tabular}{|l|l|l|l|}
\hline
\textbf{Origen} &
\textbf{Destino} &
\textbf{Protocolo} &
\textbf{Esperado} \\
\hline
PCI & PCE & ICMP & Permitido \\
\hline
PCE & PCI & ICMP & Denegado \\
\hline
PCE & Web Server & ICMP & Denegado \\
\hline
Web Server & DB Server & ICMP & Permitido \\
\hline
PCI & DB Server & ICMP & Denegado \\
\hline
Web Server & DB Server & TCP 3306 & Permitido \\
\hline
PCI & DB Server & TCP 3306 & Denegado \\
\hline
PCI & Web Server & HTTP & Permitido \\
\hline
PCE & Web Server & HTTP & Permitido \\
\hline
\end{tabular}
\caption{Matriz de resultados esperados
después de aplicar las ACL.}
\end{table}

La restricción PCE--DB Server
necesita una verificación
independiente del firewall
del servidor, al tratarse
de tráfico dentro de la
misma red local.

\subsection{Pruebas ICMP}

Para comprobar el filtrado
se utilizaron solicitudes
de ping desde diferentes
equipos.

Una prueba especialmente
importante consiste en
comparar las siguientes
comunicaciones:

\begin{itemize}
    \item PCI hacia PCE:
    debe estar permitida.

    \item PCE hacia PCI:
    debe estar denegada.
\end{itemize}

Esta diferencia demuestra
que la ACL distingue
el sentido de la solicitud ICMP.

También deben comprobarse
las solicitudes dirigidas
al servidor de bases de datos.

\subsection{Pruebas del servicio MySQL}

Desde Web Server se utilizó
el cliente MySQL/MariaDB
para establecer una conexión
con DB Server:

\begin{lstlisting}
mariadb \
    -h 140.10.20.10 \
    -u webuser \
    -p \
    laboratorio_acl
\end{lstlisting}

La consulta de los productos
permite comprobar que las reglas
de seguridad no impiden
el servicio autorizado.

\begin{lstlisting}[language=SQL]
SELECT * FROM productos;
\end{lstlisting}

\subsection{Pruebas del servicio HTTP}

La aplicación web se verificó
mediante:

\begin{lstlisting}
curl -i http://localhost/
\end{lstlisting}

El funcionamiento correcto
de la página confirma
la ejecución de PHP
y la recuperación de los
registros de la base de datos.

El acceso desde otras estaciones
se realizó mediante la dirección:

\begin{lstlisting}
http://52.20.30.10/
\end{lstlisting}

El servicio HTTP
utiliza un flujo distinto
del utilizado para la conexión SQL.

Por ello, las políticas
pueden permitir el acceso
a la página web sin autorizar
una conexión directa a
la base de datos.

\section{Inserción de un registro de prueba}

Como verificación adicional
de la integración entre
la aplicación web y
la base de datos,
se preparó una instrucción
SQL para registrar un
nuevo producto.

Dentro de MySQL/MariaDB
se seleccionó la base:

\begin{lstlisting}[language=SQL]
USE laboratorio_acl;
\end{lstlisting}

La instrucción de inserción
utilizada como prueba fue:

\begin{lstlisting}[language=SQL]
INSERT INTO productos
(nombre, precio, cantidad, fecha_registro, disponible)
VALUES
('Laptop Dell', 15999.99, 10, CURDATE(), 1);
\end{lstlisting}

Esta instrucción supone
que el campo \texttt{id}
se genera automáticamente.

Si el campo no utiliza
\texttt{AUTO\_INCREMENT},
debe proporcionarse
un identificador único.

La consulta posterior fue:

\begin{lstlisting}[language=SQL]
SELECT * FROM productos;
\end{lstlisting}

Al actualizar la página web,
el registro nuevo debe aparecer
en la tabla sin necesidad
de modificar el código PHP
ni reiniciar Apache.

Esta prueba permite comprobar
que la información presentada
es obtenida dinámicamente
desde la base de datos
y no está escrita directamente
en el documento HTML.

\captura{insertProducto.png}{
    Registro de un nuevo producto
    mediante una instrucción SQL.
}

\section{Análisis de resultados}

La implementación permitió integrar
servicios de aplicación con
mecanismos de enrutamiento
y control de acceso.

En la etapa inicial se comprobó
el funcionamiento de Apache
y PHP mediante una solicitud
HTTP local.

Posteriormente, se estableció
comunicación entre Web Server
y DB Server utilizando
MySQL/MariaDB.

La aplicación mostró correctamente
los registros obtenidos de
la tabla \texttt{productos}.

Al trasladar ambos servidores
a la topología física,
se mantuvo el funcionamiento
de la comunicación entre
aplicación y base de datos.

En cuanto al filtrado,
se configuraron dos ACL
extendidas sobre RE.

La primera establece
un bloqueo específico
de solicitudes ICMP
originadas desde la
red externa.

La segunda restringe
el tráfico enrutado
hacia DB Server,
autorizando las comunicaciones
procedentes de Web Server.

La correcta evaluación
de estas políticas requiere
comparar los resultados
de las pruebas anteriores
y posteriores a la configuración.

Los contadores de coincidencias
de las ACL constituyen
evidencia adicional para
determinar si las reglas
están procesando el tráfico.

Debe resaltarse que
el comportamiento de los
equipos dentro de una misma
subred no puede controlarse
mediante las ACL de RE
si sus paquetes no atraviesan
dicho router.

Por lo tanto, el cumplimiento
total de la restricción
de acceso a DB Server
requiere revisar también
las reglas de seguridad
del sistema operativo
o modificar la segmentación
de la red.

\section{Conclusiones}

Durante esta práctica
se integraron conocimientos
de direccionamiento IPv4,
enrutamiento dinámico,
servicios web,
bases de datos y
seguridad de redes.

La instalación de Apache
y PHP permitió implementar
una aplicación web funcional
capaz de consultar información
almacenada en un servidor remoto.

La utilización de PDO
facilitó la conexión
con MySQL/MariaDB,
mientras que la separación
de las credenciales
en un archivo externo
mejoró la organización
de la aplicación.

La implementación de OSPF
permitió comunicar
las redes interna y externa,
proporcionando una base
para evaluar el comportamiento
de las listas de control
de acceso.

Las ACL extendidas
permitieron establecer
reglas específicas
considerando direcciones
de origen y destino,
protocolos y tipos
de mensajes ICMP.

Una de las principales
observaciones fue
la importancia del sentido
de aplicación de una ACL.

Una misma interfaz
puede aplicar reglas
diferentes al tráfico
entrante y saliente,
permitiendo adaptar
el filtrado a los
requisitos de seguridad.

También se identificó
una limitación importante:
los paquetes intercambiados
entre dispositivos
de una misma VLAN
no necesariamente atraviesan
el router donde están
configuradas las ACL.

Por ello, el diseño
de políticas de seguridad
debe considerar no solamente
las direcciones IP,
sino también la topología
y los recorridos reales
del tráfico.

Finalmente, la práctica
demostró que es posible
combinar mecanismos
de filtrado en routers
con servicios de aplicación,
manteniendo las comunicaciones
autorizadas y restringiendo
aquellas que no forman parte
de las políticas establecidas.

\section{Evidencias}

\captura{evidencias1.png}{
Evidencia 1
}
\captura{evidencias2.png}{
Evidencia 2
}
\captura{evidencias3.png}{
Evidencia 3
}
\captura{evidencias4.png}{
Evidencia 4
}
\end{document}
